\documentclass[12pt]{article}
\usepackage{styles/cwilliams-standard}

% Set up document metadata
\setauthors{Tzu-Chen Chiu, Arjun Sivakumar, Mara Schimmel, Christopher Williams}
\setheaderauthors{Chiu, Sivakumar, Schimmel, Williams}
\settitle{Sprint 3 Report}
\setclass{IPD}

\begin{document}

\maketitlepage

\begin{important}{Sprint 3 Progress}
    This sprint was not as progressive in terms of getting our sprint tasks done, as we were both 
    working on completeing Finals, assignments, and projects as well as doing research for our project here. 
    We were looking into Transformers, embedding schemes for our data, and more understanding of quantum 
    computers and codes. Thank you for your understanding. 
\end{important}

\section{Original Sprint 3 Backlog}

\begin{figure}[H]
    \centering
    \includegraphics[width=1\textwidth]{images/sprint3_backlog.png}
    \caption{Spring 3 Backlog}
\end{figure}

\section{Sprint 3 Backlog - End of Sprint Status}

\begin{figure}[H]
    \centering
    \includegraphics[width=1\textwidth]{images/sprint_3_updated.png}
    \caption{Spring 3 Backlog - Updated}
    \label{fig:sprint3updated}
\end{figure}

\section{Project Management Screenshots}

The Sprint backlogs in Google Sheets are being used to manage tasks.

\section{Sprint Retrospective}

\subsection*{Mara}

As this sprint was during the end of the semester and finals week, I was unable to dedicate as much 
time as I would have preferred to this project and wasn't able to make progress during this sprint. Next 
sprint I will have a renewed focus and drive to make significant progress on this project.

\subsection*{Chris}

This sprint was particularly lack luster in terms of progress and overall accomplishment because of 
many deadlines and responsibilites in other classes. I am aware that this is not an acceptable mindset 
or outcome, especially for working life. The reason I made the decision to tackle my sprint objectives 
next sprint is in the fact that there is another sprint. The design of this project allows for us to 
shift goals and deliverables around and the extension of this class track into two semesters presents an 
opportunity like I had this semester to prioritize shorter term requirements and to look long term at what 
I will have time to accomplish next semester.

\subsection*{Tzu-Chen}

We all recently gained access to ARC with Dr. Cameron’s assistance, 
and we truly appreciate the support. Our team is currently 
learning how to connect and use ARC effectively. There are many training 
videos we are working through, including how to connect to the university 
VPN, log in via SSH, and choose appropriate compute nodes so that we do 
not overwhelm shared resources. We also appreciate Arjun for teaching 
us the basics of using ARC.

We have been developing a Transformer model based on Google’s design, 
since the one we built is not working at all. I spent substantial time 
studying how to implement the Lion optimizer, and now that I can run jobs 
on ARC, I am able to obtain logical error rate (LER) results. We hope 
to achieve improved performance in the next sprint.


\subsection*{Arjun}

During this sprint, we successfully integrated the major components developed across previous iterations, 
including the MWPM decoder, Stim data simulation, the new embedding scheme, and our transformer model, 
into a unified end-to-end pipeline. This system now runs both MWPM and transformer-based decoding on 
simulated Stim datasets, applies the analog-conversion stage, computes logical error rates, and 
generates performance graphs across varying physical error rates and code distances. We deployed and 
executed the full workflow on the ARC cluster, producing the finalized results for the team. The 
primary impediments involved merging independently developed modules into one cohesive pipeline and 
coordinating with teammates to understand how their code sections were intended to function. 
These were addressed through regular communication and reviewing each member’s demos. An 
additional challenge was that Arjun was initially the only team member with ARC access, 
requiring extra effort to run jobs and later help others onboard. From an ethical standpoint, 
we ensured proper attribution when incorporating teammates’ code into the integrated pipeline. 
Overall, the sprint went well, though future work could benefit from earlier cross-module 
synchronization and shared access to compute resources.


\section{Product Backlog (Updated)}

\begin{figure}[H]
    \centering
    \includegraphics[width=1\textwidth]{images/sprint3_updated_product_backlog.png}
    \caption{Updated Product Backlog}
\end{figure}

\begin{aside}{Picture Quality}
    We understand that this image is hard to read, but as Figure \ref{fig:sprint3updated} shows, tasks 
    1.11 was worked on.
\end{aside}

\section{Issue Tracking}

\begin{table}[H]
\centering
\begin{tabular}{@{}llll@{}}
\toprule
\textbf{Issue} & \textbf{Description} & \textbf{Status} & \textbf{Date Completed} \\ \midrule
ARC Access & Getting every member access to the ARC & 12/9/2025 \\
Train on Stim Data & Getting reasonable results in training & Ongoing \\
\end{tabular}
\caption{Issue Tracking Summary}
\end{table}

\section{Sprint 4 Backlog}

\begin{figure}[H]
    \centering
    \includegraphics[width=1\textwidth]{images/sprint_4_backlog.png}
    \caption{Spring 4 Backlog}
\end{figure}

\section{Member Videos}

Members that did not have progress for this week do not have demo videos to submit. Otherwise, 
videos will be posted on Canvas in the My Media section.

\end{document}